\section{Hyperparameters and Training Details}\label{sec:exp:hparams}

This section collects numerical training choices deferred from \cref{sec:method:training,sec:method:loss,sec:method:arch}.
Unless \cref{sec:exp:ablation} states otherwise, baseline models use the same epoch budget and hardware (\cref{sec:exp:setup}) so that differences reflect architecture and objective rather than compute.

The contrastive temperature is fixed at $\tau = \num{0.5}$ as in \cref{sec:method:loss}.
\cref{tab:exp:hparams_main} summarises the optimisation schedule and regularisation for the proposed dual-branch encoder; entries marked with an em dash will be filled before the final thesis submission.

\begin{table}[t]
  \centering
  \caption{Primary optimisation hyperparameters for the proposed model.}
  \label{tab:exp:hparams_main}
  \begin{tabular}{@{}p{0.40\linewidth}p{0.48\linewidth}@{}}
    \toprule
    Setting & Value \\
    \midrule
    Optimiser & AdamW \parencite{loshchilov2019adamw} with PyTorch defaults $(\beta_1,\beta_2) = (0.9,0.999)$ and $\epsilon = \num{e-8}$ \\
    Peak learning rate $\eta_{\max}$ & --- \\
    Warmup & Linear from $0$ to $\eta_{\max}$ over $5\%$ of total steps \\
    Schedule after warmup & Cosine annealing to $\eta_{\min}$ \\
    Minimum learning rate $\eta_{\min}$ & \num{0} \\
    Epochs & \num{8} \\
    Dropout (attention and feed-forward) & \num{0.1} throughout the transformer blocks \\
    Batch size $B$ (windows per step) & --- \\
    Mode loss weight $\lambda_{\mathrm{md}}$ & \num{1} \\
    Weight decay & --- \\
    \bottomrule
  \end{tabular}
\end{table}

Total optimisation steps equal eight epochs times the number of training batches per epoch.
The warmup fraction and cosine phase are counted in those steps, as stated in \cref{sec:method:training}.

\Cref{tab:exp:hparams_attn} collects the attention-variant knobs deferred from \cref{sec:method:attention} and the \gls{dbscan} radii of \cref{sec:method:arch}.
Vanilla and Bias do not use rotary scales.
RoPE-TOA uses a single scale $\gamma_t$ on every layer; RoPE-az uses $\gamma_t$ on the trunk and the mode branch and $(\gamma_c,\gamma_s)$ on the deinterleaving branch.
Unused coordinates are marked n/a; an em dash in an $\varepsilon$ column, or in a $\gamma$ column that the model does use, will be replaced after selection.

\begin{table}[tbp]
  \centering
  \caption{Rotary scales and \gls{dbscan} radii. $\varepsilon_{\mathrm{em}}$ and $\varepsilon_{\mathrm{md}}$ are selected independently per model on the validation set; $\mathrm{minPts}=5$ is fixed.}
  \label{tab:exp:hparams_attn}
  \begin{tabular}{@{}lccccc@{}}
    \toprule
    Model & $\gamma_t$ & $\gamma_c$ & $\gamma_s$ & $\varepsilon_{\mathrm{em}}$ & $\varepsilon_{\mathrm{md}}$ \\
    \midrule
    Vanilla   & {n/a} & {n/a} & {n/a} & {---} & {---} \\
    RoPE-TOA  & {---} & {n/a} & {n/a} & {---} & {---} \\
    RoPE-az   & {---} & {---} & {---} & {---} & {---} \\
    Bias      & {n/a} & {n/a} & {n/a} & {---} & {---} \\
    \bottomrule
  \end{tabular}
\end{table}

The single-task runs of \cref{sec:exp:ablation} inherit the Vanilla row of \cref{tab:exp:hparams_main,tab:exp:hparams_attn}, including a re-selected $\varepsilon$ per branch on the validation set of each objective.
